\begin{lemma}
Let $m\in\mathbb{N}$, let $A\in\mathbb{R}^{m\times m}$ be symmetric, i.e.\ $A^T=A$, and let
$B\in\mathbb{R}^{m\times m}$ be arbitrary. Then
$$\operatorname{Tr}\Big(A\cdot\tfrac12\big(B+B^T\big)\Big)=\operatorname{Tr}(AB).$$
Equivalently, replacing the second factor by its symmetrization does not change its trace against
a symmetric matrix. This is the matrix fact underlying the paper's remark, just after
\texttt{eq:maineq}, that because $a(\xi)$ is symmetric the equation
$\mathcal{L}u=-\operatorname{Tr}\big(a(\mathfrak{X}u)\mathfrak{X}\mathfrak{X}u\big)=0$ can also be
written with the symmetrized sub-elliptic second derivative
$\mathfrak{X}\mathfrak{X}^\star u$ in place of $\mathfrak{X}\mathfrak{X}u$.
\end{lemma}
\begin{proof}
First, $\operatorname{Tr}\big(AB^T\big)=\operatorname{Tr}(AB)$. Indeed, by the entrywise formula
for the trace of a product (preamble) and the symmetry of $A$,
$$\operatorname{Tr}\big(AB^T\big)=\sum_{i,k=1}^mA_{i,k}\big(B^T\big)_{k,i}
=\sum_{i,k=1}^mA_{i,k}B_{i,k}
=\sum_{i,k=1}^mA_{k,i}B_{i,k}=\operatorname{Tr}(AB),$$
where the first and last equalities are that formula applied to $AB^T$ and to $AB$, the second
uses $\big(B^T\big)_{k,i}=B_{i,k}$, and the third renames the summation indices
$i\leftrightarrow k$ and uses $A_{i,k}=A_{k,i}$, which is the hypothesis $A^T=A$.

Hence, using the linearity of the trace in the sense of the identities
$\operatorname{Tr}\big(A(B+C)\big)=\operatorname{Tr}(AB)+\operatorname{Tr}(AC)$ and
$\operatorname{Tr}\big(A(cB)\big)=c\operatorname{Tr}(AB)$ (preamble), with $C:=B^T$ and
$c:=\tfrac12$,
$$\operatorname{Tr}\Big(A\cdot\tfrac12\big(B+B^T\big)\Big)
=\tfrac12\operatorname{Tr}\big(A(B+B^T)\big)
=\tfrac12\operatorname{Tr}(AB)+\tfrac12\operatorname{Tr}\big(AB^T\big)
=\tfrac12\operatorname{Tr}(AB)+\tfrac12\operatorname{Tr}(AB)=\operatorname{Tr}(AB),$$
the third equality by the first paragraph.
\end{proof}
